% ==========================================================================
% Surrogate models
% ==========================================================================

\section{Surrogate Models}
\label{sec:surrogates}

The objective of this section is to investigate whether low-dimensional
surrogate models can reproduce selected effects of complex spatial
radiotherapy dynamics after PDE-to-ODE reduction and loss of spatial
information.
We compare three surrogate formulations of increasing complexity:
(i) a mechanistic two-state ODE calibrated through data assimilation,
(ii) a purely data-driven Neural ODE (NODE), and
(iii) a physics-informed residual PI-NODE combining mechanistic
radiotherapy dynamics with a learnable neural closure term.
The comparison focuses on how these approaches capture unresolved
effects associated with tumor--beam mismatch, incomplete irradiation
overlap, and delayed post-treatment dynamics.

To investigate robustness under observational uncertainty, additional
training ensembles were generated by perturbing the sparse observations
with Gaussian noise. Independent surrogate realizations were then
trained for each noisy observation set, producing probabilistic bundles
of trajectories and uncertainty bands for ODE, NODE, and PI-NODE models.

\subsection{Data Assimilation to a Reduced Two-State ODE Model}

As the first surrogate modeling strategy, we consider a reduced-order
two-state ordinary differential equation (ODE) system describing the
evolution of the normalized tumor mass together with an auxiliary latent
radiation-damage variable. The purpose of this model is to provide a
minimal mechanistic approximation of the spatial radiotherapy dynamics
generated by the reference PDE system while preserving delayed treatment
effects and multi-fraction memory.

The reduced surrogate is formulated as:
\begin{equation}
\label{eq:ode_rt_2state}
\dot y=\rho y(1-y/K_{\mathrm{eff}})-\gamma zy-\mu y,
\quad
\dot z=-z/\tau+U(t).
\end{equation}
where: \(z(t)\) is a latent accumulated radiation-damage variable, $\gamma$ is radiation-induced decay coefficient
\(\mu\) represents background decay,
and \(\tau\) is the characteristic damage relaxation time.

Unlike the earlier single-state formulation, the introduction of the
latent variable \(z(t)\) enables the surrogate to reproduce delayed
responses and cumulative effects associated with repeated irradiation
fractions. In particular, the second equation acts as a low-dimensional
memory mechanism integrating previous dose delivery.
The surrogate parameters
$\theta =
(\rho,\gamma,K,\mu,\tau)$
are estimated through data assimilation \cite{heaney2024dataassimilation} using synthetic observations
generated by the spatial PDE reference model. However, the reduced surrogate parameters should not be interpreted as physically identifiable quantities, but rather as effective closure-compensating coefficients. This affects all surrogates described in this paper. The assimilation objective is formulated as a nonlinear least-squares
optimization problem:
\begin{equation}
\label{eq:ode_training}
\theta^{*}
=
\arg\min_{\theta}
\mathcal{L}_{ODE}(\theta),
\qquad
\mathcal{L}_{ODE}(\theta)
=
\sum_{i=1}^{N_t}
\left(
y_{ODE}(t_i;\theta)
-
y_{PDE}(t_i)
\right)^2.
\end{equation}
Here, \(y_{PDE}(t_i)\) denotes the tumor mass extracted from the spatial
PDE reference simulation, whereas
\(y_{ODE}(t_i;\theta)\) is the corresponding reduced-order surrogate
prediction.
The optimization was performed using a hybrid strategy combining global
differential-evolution initialization followed by bounded local
gradient-based refinement. Assimilation was carried out only on sparse
synthetic observations sampled within a restricted temporal observation
window, while the remaining time interval was used for forecasting and
generalization assessment.

Due to the low dimensionality of the mechanistic surrogate, standard
nonlinear least-squares assimilation was sufficient in the present
study. For larger mechanistic systems with higher-dimensional parameter
spaces, more advanced assimilation approaches such as adjoint-based,
Bayesian, or ensemble methods may become necessary
\cite{heaney2024dataassimilation}.
\subsection{Classical Neural ODE (NODE) Surrogate}

While reduced mechanistic ODE surrogates provide an interpretable
approximation of the spatial tumor dynamics, their predictive capability
remains limited by the loss of spatial information during PDE-to-ODE
reduction. In particular, unresolved effects associated with
tumor--beam geometry mismatch, asymmetric irradiation, delayed treatment
response, and nonlocal spatial interactions cannot be represented
explicitly within fixed low-dimensional mechanistic models.

In the present work, the NODE surrogate \cite{chen2018neural}, \cite{rackauckas2020ude}, \cite{rubanova2019latent} is formulated as a two-state
continuous-time dynamical system:
\begin{equation}
\label{eq:node}
\dot y=f_{\psi}^{(y)}(y,z,t,U),
\qquad
\dot z=f_{\psi}^{(z)}(y,z,t,U).
\end{equation}
where:
\(f_{\psi}=(f_{\psi}^{(y)},f_{\psi}^{(z)})\) denotes a neural network
parameterized by trainable weights \(\psi\). The second latent state represents unresolved treatment-induced damage dynamics.
The latent state $z(t)$ provides low-dimensional treatment memory and delayed-response dynamics.

The neural network acts as a continuous-time nonlinear dynamical
operator approximating the temporal evolution of the reduced latent
state. In the present study, a fully connected feed-forward
architecture was employed consisting of:
\begin{itemize}
\item
an input layer containing $(y,z,t,U(t))$,
\item
two hidden layers with ReLU activation functions,
\item
a two-dimensional output representing the time derivatives $(dy/dt,dz/dt)$.
\end{itemize}
The final network layer was initialized close to zero in order to
stabilize early-stage trajectory integration and avoid unrealistically
large initial vector fields. Trajectory prediction was obtained through explicit time integration of
the learned neural dynamical system using a fourth-order Runge--Kutta
scheme. The NODE parameters were estimated by minimizing the trajectory mismatch
with respect to the PDE-generated observations:
\begin{equation}
\label{eq:node_training}
\psi^{*}
=
\arg\min_{\psi}
\mathcal{L}_{NODE}(\psi),
\qquad
\mathcal{L}_{NODE}
=
\sum_{i=1}^{N_t}
\left(
y_{NODE}(t_i)
-
y_{PDE}(t_i)
\right)^2.
\end{equation}
Only sparse observations from a restricted temporal assimilation window
were used during training, whereas the remaining time interval served as
an out-of-sample forecasting region for evaluating extrapolation
capability.
Although NODE surrogates offer substantially higher flexibility than
fixed mechanistic reduced-order models, they do not explicitly preserve
physical interpretability and may exhibit unstable extrapolation
behavior outside the assimilation interval, particularly under sparse
training conditions.

\subsection{Physics-Informed Neural ODE (PI-NODE) Surrogate}

To combine mechanistic interpretability with the flexibility of neural
correction terms, we introduce a physics-informed residual Neural ODE
(PI-NODE) surrogate. The proposed formulation augments the reduced
two-state mechanistic radiotherapy model with a learnable neural
closure term accounting for unresolved effects eliminated during the
PDE-to-ODE reduction process.

The PI-NODE surrogate augments the mechanistic reduced radiotherapy
model with a learnable neural residual closure term:
\begin{equation}
\label{eq:pinode_weighted}
\dot y
=
\omega f^{(y)}_{\mathrm{RT}}(y,z,U;\theta)
+
s_r g_{\psi}(y,z,t,U),
\qquad
\dot z
=
f^{(z)}_{\mathrm{RT}}(z,U;\theta),
\end{equation}
where
\(f^{(y)}_{\mathrm{RT}}\)
and
\(f^{(z)}_{\mathrm{RT}}\)
denote the mechanistic tumor-growth and latent damage dynamics defined
by Eq.~\eqref{eq:ode_rt_2state}.
Here:
\(\theta=(\rho,\gamma,K,\mu,\tau)\)
contains the mechanistic model parameters,
\(g_{\psi}\) denotes the neural residual closure term with trainable
parameters \(\psi\),
\(\omega\) controls the contribution of the mechanistic tumor dynamics,
and \(s_r\) scales the neural residual correction.

The residual neural component can be interpreted as a learned
reduced-order UDE closure model compensating for unresolved spatial
dynamics eliminated during PDE-to-ODE reduction, including:
incomplete overlap between tumor and irradiation regions,
delayed recolonization of irradiated areas,
outward expansion of untreated tumor regions,
nonlocal spatial effects,
and unresolved geometric heterogeneity.

In the present implementation, the neural residual correction is applied
only to the tumor evolution equation, whereas the latent damage
dynamics remain fully mechanistic.
Unlike the pure NODE formulation, the PI-NODE framework preserves
physically interpretable mechanistic parameters while allowing
unresolved dynamics to be reconstructed through residual learning.

The residual network employs the same feed-forward architecture as the
pure NODE model:
input variables $(y,z,t,U(t))$,
two hidden fully connected layers with ReLU activation functions,
and a residual output correction applied to the tumor evolution
equation.
To stabilize learning and prevent the residual correction from
overwhelming the mechanistic dynamics, the residual output was
initialized close to zero and additional \(L_2\) regularization was
imposed on the neural vector field.
\begin{equation}
\label{eq:pinode_training}
\begin{aligned}
(\theta^{*},\psi^{*})
&=
\arg\min_{\theta,\psi}
\mathcal{L}_{\mathrm{PINODE}},
\\
\mathcal{L}_{\mathrm{PINODE}}
&=
\sum_{i=1}^{N_t}
\left(
y_{\mathrm{PINODE}}(t_i)
-
y_{PDE}(t_i)
\right)^2
+
\lambda\,\mathcal{R}(g_{\psi}).
\end{aligned}
\end{equation}
where $\mathcal{R}(g_{\psi})
=\|g_{\psi}\|_2^2$
denotes residual regularization preventing the neural correction from
dominating the mechanistic component.
Training was performed using an alternating optimization strategy. In
particular, residual-network parameters and mechanistic parameters were
updated in separate training phases followed by a final joint
fine-tuning stage. Mechanistic parameters were initialized directly from
the reduced PDE-derived physical model rather than from the assimilated
ODE fit.

The implementation strategy was inspired by the Scientific Machine
Learning (SciML) ecosystem and related Universal Differential Equation
(UDE) \cite{rackauckas2019diffeqflux}, \cite{rackauckas2020ude} formulations combining mechanistic differential equations with
trainable neural closure operators.
